\section{おわりに}
\label{sec:conclusion}

同期固定カメラからテニスのプレイヤーとボールを同じコート座標に復元する標準経路を示した。
規格コートで視点を整列し、選手を対応付けて関節と球を三角測量する。身体姿勢、位置、球軌道を元動画の時間軸で統合し、
復元の有効性と失敗理由を保存する。
コートの学習シーン構成、選手選別、球検出の時間文脈を比較し、下流の観測を左右する要因を示した。
宣言型計画と成果物保存でラリー単位の処理を動画全体へ展開し、モデル共有とTensorRTによる
時間・資源の変化を、欠測と座標差の実測とともに報告した。
